% !TeX program = tectonic
\documentclass[11pt,a4paper]{article}
\usepackage{fontspec}
\usepackage[english]{babel}
\babelfont{rm}[Language=Default]{Liberation Serif}
\babelfont[Language=Default]{sf}{Carlito}
\babelfont[Language=Default]{tt}{DejaVu Sans Mono}
\usepackage{amsmath,amssymb,amsthm}
\usepackage{geometry}
\geometry{a4paper, top=2.4cm, bottom=2.4cm, left=2.4cm, right=2.4cm}
\usepackage{booktabs,tabularx,enumitem,xcolor,tcolorbox}
\tcbuselibrary{breakable,skins}
\definecolor{linkblue}{HTML}{1F4E79}
\definecolor{statgreen}{HTML}{2F6B3A}
\definecolor{goldacc}{HTML}{8B7E5A}
\usepackage[colorlinks=true,linkcolor=linkblue,urlcolor=linkblue,unicode=true]{hyperref}
\theoremstyle{plain}
\newtheorem{theorem}{Theorem}
\newtheorem{lemma}[theorem]{Lemma}
\newtheorem{proposition}[theorem]{Proposition}
\newtheorem{corollary}[theorem]{Corollary}
\theoremstyle{definition}
\newtheorem{definition}[theorem]{Definition}
\theoremstyle{remark}
\newtheorem{remark}[theorem]{Remark}
\newtcolorbox{statusbox}[1][]{colback=statgreen!5,colframe=statgreen!75,
  title={\textbf{Summary of results:} #1},fonttitle=\small,boxrule=0.7pt,arc=2pt,breakable}
\newtcolorbox{databox}[1][]{colback=goldacc!6,colframe=goldacc!80,
  title={\textbf{Protocol data:} #1},fonttitle=\small,boxrule=0.7pt,arc=2pt,breakable}
\newtcolorbox{verifybox}[1][]{colback=linkblue!5,colframe=linkblue!80,
  title={\textbf{Verification:} #1},fonttitle=\small,boxrule=0.7pt,arc=2pt,breakable}
\widowpenalty=10000
\clubpenalty=10000
\setlength{\parskip}{0.35em}
\newcommand{\CC}{\mathbb{C}}
\newcommand{\RR}{\mathbb{R}}
\newcommand{\ZZ}{\mathbb{Z}}
\newcommand{\QQ}{\mathbb{Q}}
\newcommand{\Ga}{\Gamma}
\newcommand{\Bfun}{\mathrm{B}}
\newcommand{\im}{\operatorname{Im}}
\newcommand{\re}{\operatorname{Re}}
\newcommand{\lcm}{\operatorname{lcm}}
\newcommand{\SNF}{\operatorname{SNF}}
\newcommand{\Jac}{\operatorname{Jac}}
\newcommand{\rank}{\operatorname{rank}}
\newcommand{\Ind}{\operatorname{Index}}
\newcommand{\disc}{\operatorname{disc}}
\begin{document}
\begin{center}
{\LARGE\bfseries The Genus of the Fermat Curve}\\[0.4em]
{\large the topological foundation of the cyclotomic stands}\\[0.6em]
{\normalsize Isaev Iskhak Khamzatovich}\\[0.2em]
{\normalsize The <<Dynamic Principle>> Program --- the \texttt{hodge-laboratory}}\\[0.15em]
{\small Standalone theorem monograph · Монография 1/16}
\end{center}
\vspace{0.4em}
{\color{linkblue}\hrule height 0.8pt}
\vspace{0.8em}

\section{Setting}

The Fermat curve $C_N=\{x^N+y^N=z^N\}\subset\mathbb{P}^2$ is the
basic object of the cyclotomic stands. Its topology is determined by
a single integer $N$, and the first pipeline step --- the genus $g$
--- is purely topological: no analysis, no spectral theory. The
genus determines the dimension of the space of holomorphic forms,
the size of the period matrix, and ultimately the scale of the whole
construction.

\begin{lemma}[genus of the Fermat curve]\label{lem:genus}
Let $N\ge 4$. The projection $\pi_x\colon C_N\to\mathbb{P}^1$ is a
morphism of degree $N$, unramified over $\infty$ and ramified
exactly over the $N$ points $\zeta_N^k$, with ramification index $N$
at each $P_k=(\zeta_N^k,0)$. The genus equals
$g=\frac{(N-1)(N-2)}{2}$; in particular $g(15)=91$, $g(30)=406$.
\end{lemma}

\begin{proof}
Over any $x_0\in\mathbb{P}^1$ the equation $y^N=1-x_0^N$ has exactly
$N$ solutions with multiplicity, hence $\deg\pi_x=N$. At $P_k$ the
local expansion $1-x^N=-N\zeta_N^{k(N-1)}(x-\zeta_N^k)+\dots$ gives
$y^N=c\,(x-\zeta_N^k)$ with $c\ne0$: the parameter is $y$ and
$x-\zeta_N^k\sim c^{-1}y^N$, so the ramification index is $N$. Over
$\infty$ there is no ramification: at $z=0$ the equation
$X^N+Y^N=0$ has $N$ distinct solutions with unramified projection.
The Riemann--Hurwitz formula gives
$2g-2=-2N+N(N-1)=N^2-3N$, hence
$g=(N-1)(N-2)/2$: $91$ and $406$ at $N=15,30$.
\end{proof}

\begin{remark}
The same Pl\"ucker formula for a smooth plane curve of degree $m$
gives $\tfrac{(m-1)(m-2)}2$ --- the agreement of the two readings
serves as a built-in control: at $N=4$ both give $3$ (the Klein
quartic is also a genus-3 curve, but a \emph{different} curve).
\end{remark}

\begin{databox}[numbers]
$g(15)=91$; period matrix $91\times182$; $g(30)=406$; period matrix
$406\times812$; for comparison: $g(4)=3$, $g(7)=15$, $g(22)=231$.
\end{databox}

\begin{statusbox}[summary]
The genus is proved from Riemann--Hurwitz with all components
explicit. Consequences: the character censuses $91=1+6+84$ and
$406=1+6+9+30+84+276$ agree with the genera (check V1, exact
integer arithmetic).
\end{statusbox}

\begin{verifybox}[reproduction]
\texttt{python3 laboratory.py} $\to$ ``Stand N=15/30 $\to$ V1
census''; Lean: \texttt{genus\_fermat} in
\texttt{HodgeLaboratory.lean} --- machine proofs of
$(15-1)(15-2)/2=91$ and $(30-1)(30-2)/2=406$.
\end{verifybox}

\vfill
{\color{linkblue}\hrule height 0.6pt}
\vspace{0.5em}
{\small\color{gray}
License: individual exclusive license of Isaev Iskhak Khamzatovich.
All rights to the computations, formulas, and texts belong to the author.
Verification: \texttt{python3 laboratory.py} (``Stands''/``Certificates''),
Lean: \texttt{verification/lean/HodgeLaboratory.lean}.
}
\end{document}
